\label{sec:dhondt}

\begin{theorem}[Jefferson-D'Hondt Equivalence]
  \label{thm:dhondt-equiv}
  The Jefferson apportionment method and the d'Hondt proportional
  representation method are mathematically identical.
  For any set of populations (or vote totals) $p_1, \ldots, p_n$
  and total seats $H$, both methods produce the same seat allocation
  $(s_1, \ldots, s_n)$.
\end{theorem}

\begin{proof}
  Jefferson's priority for state $i$'s $(s+1)$-th seat is
  $P_i^{\text{Jeff}}(s) = p_i/(s+1)$.
  D'Hondt's priority for party $j$'s next seat when it has $s$ seats is
  $P_j^{\text{D'H}}(s) = V_j/(s+1)$.
  With $p_i = V_j$ (population = vote total) and $n = $ (number of
  states) $= $ (number of parties), the two priority formulas are
  identical: $P_i^{\text{Jeff}}(s) = P_j^{\text{D'H}}(s)$ for all $s$.
  Since the priority-queue algorithm depends only on the priority
  function, both methods execute identical sequences of seat awards
  and produce identical allocations.
\end{proof}

The equivalence is an instance of a more general structural identity:
any divisor method applied to a collection of ``vote shares''
(whether those shares are state populations, party vote totals,
or any other quantity to be proportionally allocated) produces
the same seat allocation.
The only formal difference between the U.S.\ apportionment context
and the PR election context is the constitutional floor:
U.S.\ apportionment requires $s_i \geq 1$ for all states, while
PR elections do not require every party to receive a seat.

\subsection{The D'Hondt Method in Practice}

D'Hondt's original 1882 paper \citep{dhondt1882systeme} proposed the
method for Belgian elections, and it was adopted by Belgium and
subsequently spread to much of Europe, Latin America, and beyond.
Countries using d'Hondt for national legislative elections include:
\begin{itemize}
  \item \textbf{Western Europe}: Belgium, Spain, Portugal, Austria,
    the Netherlands (historically), Switzerland (partially)
  \item \textbf{Eastern Europe}: Czech Republic, Poland, Romania,
    Croatia, Serbia
  \item \textbf{Latin America}: Brazil, Argentina, Colombia, Ecuador
  \item \textbf{Other}: Japan (lower house), Turkey (modified d'Hondt),
    Israel (modified)
\end{itemize}
The method is known in the literature as d'Hondt in the electoral
context and Jefferson in the U.S.\ apportionment context.
The failure to recognise the equivalence for much of the 20th century
led to redundant analysis in two separate literatures
\citep{balinski1982fair, lijphart1994electoral}.
Balinski and Young \citeyearpar{balinski1982fair} provided the
definitive unified treatment showing that the two bodies of work
address the same mathematical problem.

\subsection{Why D'Hondt Favors Large Parties}

The equivalence proof explains why d'Hondt is widely criticised
in the PR literature for favouring large parties (or, in the U.S.\
context, large states):
floor rounding discards fractional seats that would disproportionately
benefit small parties/states.
A party with 10\% of votes in a 10-seat system has an exact quota of 1.0;
a party with 10.1\% has a quota of 1.01.
Under d'Hondt, both receive 1 seat.
Under Hamilton (largest remainders), the 10.1\% party has a
larger fractional remainder (0.01 vs.\ 0.0) and is more likely to
receive a remainder seat.
Jefferson/d'Hondt's indifference to fractional remainders creates
the systematic large-party/state bias documented in Section~\ref{sec:bias}.
